\begin{lemma}
Assume the main-independence hypothesis: for every real \(\varepsilon>0\), all
sufficiently large \(n\) admit a graph \(G\) on \(\mathrm{Fin}(n)\) with
\(G.\mathrm{CliqueFree}(3)\) and
\[
\noderef{IndependenceNumberLess}
\left(G,(1+\varepsilon)\sqrt{(n:\mathbb R)\log(n:\mathbb R)}\right).
\]
Then, for every real \(0<\eta<1/2\), there is an integer \(k_0\) such that for
every integer \(k\ge k_0\) there is an integer \(N\) satisfying
\[
\noderef{RamseyWitness}(N,k)
\qquad\text{and}\qquad
\left(\frac12-\eta\right)\noderef{RamseyScale}(k)\le (N:\mathbb R).
\]
\end{lemma}
\begin{proof}
Fix \(0<\eta<1/2\), and put
\[
c=\frac12-\frac{\eta}{2}.
\]
Then
\[
\frac12-\eta<c<\frac12,\qquad
\delta:=c-\left(\frac12-\eta\right)=\frac{\eta}{2}>0.
\]
Because \(2c<1\), choose \(\varepsilon>0\) so small that
\[
2c(1+\varepsilon)^2<1. \tag{1}
\]
Apply the main-independence hypothesis to this \(\varepsilon\).  There is an
integer \(n_0\) such that every \(n\ge n_0\) has a graph \(G\) on
\(\mathrm{Fin}(n)\) with \(G.\mathrm{CliqueFree}(3)\) and
\[
\noderef{IndependenceNumberLess}
\left(G,(1+\varepsilon)\sqrt{(n:\mathbb R)\log(n:\mathbb R)}\right). \tag{2}
\]

Choose \(k_0\) large enough that, for every \(k\ge k_0\),
\[
k\ge3,\qquad \log k>0,\qquad \log k>c,
\]
\[
c\,\frac{k^2}{\log k}\ge n_0+1,\qquad
c\,\frac{k^2}{\log k}\ge4,\qquad
\delta\,\frac{k^2}{\log k}\ge1. \tag{3}
\]
Such a threshold exists because \(\log k\to\infty\) and
\(k^2/\log k\to\infty\).

Fix \(k\ge k_0\), abbreviate
\[
X=c\,\frac{k^2}{\log k},
\]
and set
\[
N=\lfloor X\rfloor. \tag{4}
\]
Since \(X\ge0\), the floor inequalities give
\[
(N:\mathbb R)\le X,\qquad X-1\le (N:\mathbb R). \tag{5}
\]
The last inequality in (3) gives
\[
1\le \delta\,\frac{k^2}{\log k}.
\]
Subtracting this from \(X=c\,k^2/\log k\), we get
\[
X-1\ge
\left(c-\delta\right)\frac{k^2}{\log k}
=\left(\frac12-\eta\right)\frac{k^2}{\log k}. \tag{6}
\]
Combining (5) and (6), and using the definition of \(\noderef{RamseyScale}\),
\[
\left(\frac12-\eta\right)\noderef{RamseyScale}(k)\le (N:\mathbb R). \tag{7}
\]
The first middle inequality in (3) also gives
\[
(N:\mathbb R)\ge X-1\ge n_0,
\]
so \(N\ge n_0\).  The second middle inequality in (3) gives
\[
(N:\mathbb R)\ge X-1\ge3,
\]
so \(N\ge3\).

We next prove that the graph supplied by (2) for this \(N\) has no independent
set of cardinality \(k\).  From \(N\ge3\), \(N>0\) and \(\log N>0\).  From
\((N:\mathbb R)\le X\) and \(\log k>c>0\),
\[
(N:\mathbb R)\le c\,\frac{k^2}{\log k}<k^2. \tag{8}
\]
Thus \(0<N\le k^2\).  Monotonicity of logarithm on positive reals gives
\[
\log N\le \log(k^2)=2\log k. \tag{9}
\]
Multiplying the upper bounds in (5) and (9), and using \(\log k>0\), yields
\[
N\log N\le
\left(c\,\frac{k^2}{\log k}\right)(2\log k)=2c\,k^2. \tag{10}
\]
Multiplying by \((1+\varepsilon)^2>0\) and using (1),
\[
(1+\varepsilon)^2N\log N<k^2. \tag{11}
\]
Both sides are nonnegative and \(k>0\), so monotonicity of square root gives
\[
(1+\varepsilon)\sqrt{N\log N}<k. \tag{12}
\]

Apply (2) at \(n=N\), using \(N\ge n_0\).  We get a graph \(G\) on
\(\mathrm{Fin}(N)\) with \(G.\mathrm{CliqueFree}(3)\) and
\[
\noderef{IndependenceNumberLess}
\left(G,(1+\varepsilon)\sqrt{N\log N}\right).
\]
If \(I\subseteq\mathrm{Fin}(N)\) were independent with \(|I|=k\), then
\(\noderef{IndependenceNumberLess}\) and (12) would give
\[
(k:\mathbb R)=(|I|:\mathbb R)
<(1+\varepsilon)\sqrt{N\log N}<k,
\]
a contradiction.  Hence no independent \(k\)-set exists in \(G\), and together
with \(G.\mathrm{CliqueFree}(3)\) this is exactly
\(\noderef{RamseyWitness}(N,k)\).  The scale lower bound is (7), so the desired
conclusion holds for every \(k\ge k_0\).
\end{proof}
